% #############################################################################
% Resumo
% !TEX root = ../../Main.tex
% #############################################################################
\acresetall
\begin{otherlanguage}{portuguese}
\noindent A aprendizagem por reforço aplicada aos mercados cambiais é normalmente testada num único par, com dados em \textit{bars} e sem custos reais. Esta tese treina um agente Deep Deterministic Policy Gradient (DDPG), com um único desenho e sem ajustes por par, em cada um dos sete principais pares de moeda, com onze anos de dados ao \textit{tick}, de 2015 a 2025. O agente é implementado como no artigo original do DDPG: um método ator-crítico \textit{off-policy} com redes do ator e do crítico, cópias-alvo de ambas atualizadas lentamente, memória de repetição e ruído de exploração de Ornstein-Uhlenbeck, com os hiperparâmetros publicados onde se aplicam. O ator transforma 38 variáveis adimensionais, obtidas de dezasseis indicadores técnicos, numa única ação contínua em $[-1, 1]$, cujo sinal define a direção da posição e cuja magnitude define a sua dimensão, dentro de um orçamento de risco baseado na volatilidade. O crítico aprende o valor dessa ação, e a recompensa é o retorno logarítmico do capital. Quatro alterações documentadas adaptam-no a este problema: redes mais pequenas com normalização por camada, uma penalização contra a saturação da saída, uma recompensa escalada e episódios de treino espelhados. O agente é treinado e avaliado num motor de \textit{backtesting} ao \textit{tick}, construído para reproduzir o do cTrader, que percorre 1,6 mil milhões de cotações e cobra os custos reais de uma conta de retalho \textit{raw-spread} e \textit{swap-free}: o \textit{spread} de cada execução, uma comissão de 3,5 dólares americanos por lote e por lado e um \textit{swap} nulo. Cada par é treinado com uma validação \textit{walk-forward} ancorada, o método padrão para sistemas de negociação: sete partições validam de 2018 a 2024, 2025 é reservado como ano de teste, e cada modelo tem de passar onze critérios de retorno e comportamento. Os sete modelos são lucrativos após todos os custos. Obtiveram entre \ResultStitchedMinPT{} e \ResultStitchedMaxPT{} nos anos de validação, entre \ResultTestMinPT{} e \ResultTestMaxPT{} no ano de teste e entre \ResultFullMinPT{} e \ResultFullMaxPT{} no período completo, com rácios de Sharpe de \ResultSharpeMinPT{} a \ResultSharpeMaxPT{} e um alfa anual de \ResultAlphaMinPT{} a \ResultAlphaMaxPT{}, e cada um supera o seu par nos três períodos. Numa carteira de pesos iguais obtiveram \ResultPortfolioReturnPT{} contra \ResultBenchmarkReturnPT{} do índice US 500, com uma perda máxima de \ResultPortfolioDrawdownPT{} contra \ResultBenchmarkDrawdownPT{} e um rácio de Sharpe de \ResultPortfolioSharpePT{} contra \ResultBenchmarkSharpePT{}, porque os sete modelos ganham e perdem dinheiro em momentos diferentes.
\end{otherlanguage}
